%%%%%%%%%%%%%%%%%%%%%%%%%%%%%%%%%%%%%%%%%%%%%%%%%%%%%%%%%%%%%%%%%%%
%%%
%%% File ini ditujukan untuk menulis 
%%% Paper POMITS menggunakan LaTeX untuk program sarjana
%%% di Institut Teknologi Sepuluh Nopember, Surabaya.
%%%
%%% Komentar, saran, koreksi, modifikasi untuk file ini dipersilakan
%%%
%%% Didapat dari: Ahmad Zaki Al Muntazhar(
%%% Matematika ITS 2019
%%% Revisi terahir oleh: Ahmad Hisbu Zakiyudin
%%% Matematika ITS 2020
%%% <azakiyudin1790@gmail.com>
%%% Revisi terakhir (Juni 2024): 
%%% 1. Penyesuaian Header (masih belum sama persis), 
%%% 2. Penyesuaian Penomoran Definisi, Teorema, Contoh, dan Persamaan (2022),
%%% 3. Penyempurnaan sesuai Format POMITS.
%%%
%%%%%%%%%%%%%%%%%%%%%%%%%%%%%%%%%%%%%%%%%%%%%%%%%%%%%%%%%%%%%%%%%%%%
\documentclass[journal]{IEEEtran}
% correct bad hyphenation here
\hyphenation{op-tical net-works semi-conduc-tor}
\usepackage{amssymb,amsmath,amsthm,graphicx}
\usepackage{caption,tabularx,multicol,multirow}
\usepackage{listings,color,booktabs}
\usepackage{float}
\usepackage{algpseudocode}
\usepackage{array}
\usepackage{scrfontsizes}
\usepackage{afterpage}
\usepackage{subfig}
\usepackage{mathtools}
\usepackage{physics}
\usepackage[ruled,vlined]{algorithm2e}
\usepackage{algpseudocode}
\captionsetup[figure]{labelsep=period}
\captionsetup[figure]{name={Gambar}}

\captionsetup[table]{labelsep=newline}
\captionsetup[table]{name={Tabel}}
\tolerance=1
\emergencystretch=\maxdimen
\hyphenpenalty=10000
\hbadness=10000
\setlength{\parindent}{0.36cm}
\newenvironment{breakablealgorithm}
  {% \begin{breakablealgorithm}
   \begin{center}
     \refstepcounter{algorithm}% New algorithm
     \hrule height.8pt depth0pt \kern2pt% \@fs@pre for \@fs@ruled
     \renewcommand{\caption}[2][\relax]{% Make a new \caption
       {\raggedright\textbf{\fname@algorithm~\thealgorithm} ##2\par}%
       \ifx\relax##1\relax % #1 is \relax
         \addcontentsline{loa}{algorithm}{\protect\numberline{\thealgorithm}##2}%
       \else % #1 is not \relax
         \addcontentsline{loa}{algorithm}{\protect\numberline{\thealgorithm}##1}%
       \fi
       \kern2pt\hrule\kern2pt
     }
  }{% \end{breakablealgorithm}
     \kern2pt\hrule\relax% \@fs@post for \@fs@ruled
   \end{center}
  }


\newtheorem{defn}{Definisi}[section]
\newtheorem{teo}[defn]{Teorema}
\newtheorem{thm}{Teorema}
\newtheorem{lemma}{Lemma}
\newtheorem{lemmas}[defn]{Lemma}
\newtheorem{cor}[defn]{Akibat}
\newtheorem{asumsi}[defn]{Asumsi}
\theoremstyle{definition}
\newtheorem{con}{Contoh}[section]
\theoremstyle{definition}
\theoremstyle{plain}
\newtheorem{prop}{Proposisi}
\renewcommand{\proofname}{Bukti}
\renewcommand{\thedefn}{\arabic{section}.\arabic{defn}}
\renewcommand{\thecon}{\arabic{section}.\arabic{con}}

\newcommand{\normlp}[1]{\left\|#1\right\|_{L^p_{\lambda}(0,T;H)}}% Fungsi norm (||x||)
\newcommand{\R}{\mathbb{R}}
\renewcommand{\natural}{\mathbb{N}}
\newcommand{\Xpl}{X^R_{p,\lambda}}
\newcommand{\Xpls}{X^R_{p,\lambda^*}}
\newcommand{\qpl}{q^R_{p,\lambda}}
\newcommand{\qpls}{q^R_{p,\lambda^*}}
\newcommand{\katakunci}[1]{%
    \noindent\\
    \textit{\textbf{Kata Kunci---}}{\textbf{#1}}
}

%\renewcommand\thetheorem{\arabic{section}.\arabic{theorem}}
\renewcommand{\listfigurename}{\hfill\normalsize\textbf{DAFTAR GAMBAR}\hfill}
\renewcommand{\listtablename}{\hfill\normalsize\textbf{DAFTAR TABEL}\hfill}
\renewcommand{\contentsname}{\hfill\normalsize\textbf{DAFTAR ISI}\hfill}
\renewcommand{\refname}{}
\renewcommand{\arraystretch}{1.2}

\newcolumntype{L}[1]{>{\raggedright\let\newline\\\arraybackslash\hspace{0pt}}m{#1}}
\newcolumntype{C}[1]{>{\centering\let\newline\\\arraybackslash\hspace{0pt}}m{#1}}
\newcolumntype{R}[1]{>{\raggedleft\let\newline\\\arraybackslash\hspace{0pt}}m{#1}}

\definecolor{codegreen}{rgb}{0,0.6,0}
\definecolor{codegray}{rgb}{0.5,0.5,0.5}
\definecolor{codepurple}{rgb}{0.58,0,0.82}
\definecolor{ITSblue}{rgb}{0.08,0.38,0.74}

\lstdefinestyle{mypseudo}{ 
commentstyle=\bfseries,
keywordstyle=\bfseries,
breakatwhitespace=false,         
breaklines=true,                 
captionpos=b,                    
keepspaces=true,                  
showspaces=false,                
showstringspaces=false,
showtabs=false,
literate={\ \ }{{\ }}2
}
\lstset{basicstyle=\sffamily\scriptsize,style=mypseudo}
\def\abstractname{Abstrak}
%\renewcommand{\thetable}{\arabic{table}}

\begin{document}

    % paper title
    % can use linebreaks \\ within to get better formatting as desired
    % Do not put math or special symbols in the title.
    \title{\changefontsizes{24}Aplikasi Teorema Titik Tetap pada Ruang Metrik Kuasi Untuk Masalah Eksistensi dan Ketunggalan Solusi dari Persamaan Cauchy Tidak Homogen\\ \vspace{12pt}}
    %
    %
    % author names and IEEE memberships
    % note positions of commas and nonbreaking spaces ( ~ ) LaTeX will not break
    % a structure at a ~ so this keeps an author's name from being broken across
    % two lines.
    % use \thanks{} to gain access to the first footnote area
    % a separate \thanks must be used for each paragraph as LaTeX2e's \thanks
    % was not built to handle multiple paragraphs
    %
    
    \author{\changefontsizes{11}Ahmad Hisbu Zakiyudin, Kistosil Fahim\\
    Departemen Matematika\\Institut Teknologi Sepuluh Nopember (ITS)\\
    Jl. Arief Rahman Hakim, Surabaya 60111 Indonesia\\
    \textit{e-mail} : azakiyudin1790@gmail.com, 	kfahim@matematika.its.ac.id\\ \vspace{-1em}}
    \markboth{\small JURNAL TEKNIK ITS Vol. X, No. Y, (TAHUN) ISSN: 2337-3539 (2301-9271 Print)}{}
    % <-this % stops a space
    %\thanks{M. Shell is with the Department
    %of Electrical and Computer Engineering, Georgia Institute of Technology, Atlanta,
    %GA, 30332 USA e-mail: (see http://www.michaelshell.org/contact.html).}% <-this % stops a space
    %\thanks{J. Doe and J. Doe are with Anonymous University.}% <-this % stops a space
    %\thanks{Manuscript received April 19, 2005; revised December 27, 2012.}}
    
    % note the % following the last \IEEEmembership and also \thanks - 
    % these prevent an unwanted space from occurring between the last author name
    % and the end of the author line. i.e., if you had this:
    % 
    % \author{....lastname \thanks{...} \thanks{...} }
    %                     ^------------^------------^----Do not want these spaces!
    %
    % a space would be appended to the last name and could cause every name on that
    % line to be shifted left slightly. This is one of those "LaTeX things". For
    % instance, "\textbf{A} \textbf{B}" will typeset as "A B" not "AB". To get
    % "AB" then you have to do: "\textbf{A}\textbf{B}"
    % \thanks is no different in this regard, so shield the last } of each \thanks
    % that ends a line with a % and do not let a space in before the next \thanks.
    % Spaces after \IEEEmembership other than the last one are OK (and needed) as
    % you are supposed to have spaces between the names. For what it is worth,
    % this is a minor point as most people would not even notice if the said evil
    % space somehow managed to creep in.
    
    
    
    % The paper headers
    %
    %{Shell \MakeLowercase{\textit{et al.}}: Bare Demo of IEEEtran.cls for Journals}
    % The only time the second header will appear is for the odd numbered pages
    % after the title page when using the twoside option.
    % 
    % *** Note that you probably will NOT want to include the author's ***
    % *** name in the headers of peer review papers.                   ***
    % You can use \ifCLASSOPTIONpeerreview for conditional compilation here if
    % you desire.
    
    
    
    
    % If you want to put a publisher's ID mark on the page you can do it like
    % this:
    %\IEEEpubid{0000--0000/00\$00.00~\copyright~2012 IEEE}
    % Remember, if you use this you must call \IEEEpubidadjcol in the second
    % column for its text to clear the IEEEpubid mark.
    
    
    
    % use for special paper notices
    %\IEEEspecialpapernotice{(Invited Paper)}
    
    
    
    
    % make the title area
     \maketitle
    % As a general rule, do not put math, special symbols or citations
    % in the abstract or keywords.
    
    \begin{abstract}
   \small \noindent Abstrak dituliskan dengan paragraf tunggal. Abstrak mencakup pendahuluan, metode dan hasil yang dicapai, tanpa ada acuan pada daftar pustaka. Abstrak harus menggambarkan penelitian yang dilakukan secara ekplisit dengan kalimat yang lugas dan jelas. Abstrak dan artikel ditulis dalam bahasa Indonesia baku.  Panjang abstrak yang disarankan adalah antara 100 hingga 300 kata.
    \end{abstract}
   % \abstrak{Abstrak dituliskan dengan paragraf tunggal. Abstrak mencakup pendahuluan, metode dan hasil yang dicapai, tanpa ada acuan pada daftar pustaka. Abstrak harus menggambarkan penelitian yang dilakukan secara ekplisit dengan kalimat yang lugas dan jelas. Abstrak dan artikel ditulis dalam bahasa Indonesia baku.  Panjang abstrak yang disarankan adalah antara 100 hingga 300 kata.}
    % Note that keywords are not normally used for peerreview papers.
    \katakunci{\small Persamaan Cauchy Tidak Homogen, Ruang Metrik Kuasi, Teorema Titik Tetap.}
   % \begin{IEEEkeywords}
   %  Quantum K-Nearest Neighbors, Komputasi Kuantum, Pengenalan Objek, Fidelity
   % \end{IEEEkeywords}

    
    
    % For peer review papers, you can put extra information on the cover
    % page as needed:
    % \ifCLASSOPTIONpeerreview
    % \begin{center} \bfseries EDICS Category: 3-BBND \end{center}
    % \fi
    %
    % For peerreview papers, this IEEEtran command inserts a page break and
    % creates the second title. It will be ignored for other modes.
    \IEEEpeerreviewmaketitle
    
     \section{PENDAHULUAN}
     \normalsize \IEEEPARstart{P}{ADA} tahun 1906, Maurice Frechet memperkenalkan ruang metrik yang menjadi titik awal perkembangan dari analisis fungsional. Seiring berjalanannya waktu, banyak matematikawan yang memperkenalkan jenis-jenis dari ruang metrik dan mengembangkannya. Salah satunya adalah ruang metrik kuasi yang diperkenalkan oleh Wilson. 

    Ruang metrik kuasi merupakan ruang metrik tanpa sifat simetri. Secara umum, semua ruang metrik merupakan ruang metrik kuasi, tetapi tidak berlaku sebaliknya \cite{Wilson1931}. Beberapa konsep pada ruang metrik kuasi sedikit berbeda dengan ruang metrik biasa. Terdapat dua konsep topologi, yaitu topologi maju dan topologi mundur, sehingga terdapat pula konsep konvergensi maju, kelengkapan maju, dan lain sebagainya \cite{JuliaCollins2007}. Akhir-akhir ini, ruang metrik kuasi memiliki banyak aplikasi pada berbagai bidang, di antaranya tentang eksistensi persamaan Hamilton-Jacobi, dalam model bebas laju untuk plastisitas, model kegagalan material, konstruksi taksonomi otomatis, dan lain-lain. \cite{Abdelkarim2023}   
    
    \section{METODE PENELITIAN}
    \hskip 0.5cm Sebelum membahas mengenai titik tetap dalam ruang metrik kuasi, berikut ini diberikan definisi dari ruang metrik kuasi beserta contohnya.
\begin{defn}\label{ruangmetrikkuasi}
    \cite{JuliaCollins2007} Diberikan himpunan tak kosong $X$ dan $q$ adalah metrik yang didefinisikan sebagai fungsi $q :X\times X\rightarrow \R$ sehingga untuk setiap $x,y,z\in X$ memenuhi 
    \begin{enumerate}
        \item[(M1)] $q(x,y)\geq 0$,
        \item[(M2)] $q(x,y)=d(y,x)=0$ jika dan hanya jika $x=y$,
        \item[(M4)] $q(x,y)\leq q(x,z)+q(z,y)$ (ketaksamaan segitiga).
    \end{enumerate}
    Pasangan terurut $(X,q)$ disebut ruang metrik kuasi.
\end{defn}
Berdasarkan definisi tersebut, berikut ini disajikan beberapa contoh dari ruang metrik kuasi
\begin{con}\label{con:ruangmetrikkuasi1}
    \cite{JuliaCollins2007} Misalkan $\alpha$ bilangan real positif. Diberikan himpunan $X:=\R$ dan untuk setiap $x,y\in X$ didefinisikan fungsi $q:X\times X\rightarrow \R$ dengan 
    \begin{align}
        q(x,y) := \begin{cases}
        y-x, &\text{ jika } y\geq x,\\
        \alpha(x-y), &\text{ jika } y<x.
        \end{cases} \label{fungsicon:ruangmetrikkuasi1}
    \end{align}
    Pasangan terurut $(X,q)$ merupakan ruang metrik kuasi.
\end{con}
\begin{con}\label{con:ruangmetrikkuasi2}
    Misalkan $\alpha\geq 0$ dan $f:\R\rightarrow\R$ fungsi naik tegas. Diberikan himpunan $X:=\R$ dan untuk setiap $x,y\in X$ didefinisikan fungsi $q:X\times X\rightarrow \R$ dengan 
    \begin{align}
        q(x,y) := \begin{cases}
        f(y)-f(x), &\text{ jika } y\geq x,\\
        \alpha(f(x)-f(y)), &\text{ jika } y<x.
        \end{cases} \label{fungsicon:ruangmetrikkuasi2}
    \end{align}
    Pasangan terurut $(X,q)$ merupakan ruang metrik kuasi.
\end{con}
\begin{con}\label{con:ruangmetrikkuasi3}
    \cite{JuliaCollins2007} Diberikan himpunan $X:=\R$ dan untuk setiap $x,y\in X$ didefinisikan fungsi $q:X\times X\rightarrow \R$ dengan 
    \begin{align}
        q(x,y) := \begin{cases}
        y-x, &\text{ jika } y\geq x,\\
        1, &\text{ jika } y<x.
        \end{cases} \label{fungsicon:ruangmetrikkuasi3}
    \end{align}
    Pasangan terurut $(X,q)$ merupakan ruang metrik kuasi.
\end{con}
Pada beberapa contoh yang telah diberikan, yaitu Contoh \ref{con:ruangmetrikkuasi1}, \ref{con:ruangmetrikkuasi2}, dan \ref{con:ruangmetrikkuasi3}, jelas bahwa $q(x,y)\neq q(y,x)$ sehingga ketiga contoh tersebut merupakan ruang metrik kuasi yang bukan ruang metrik. Dalam hal ini, Contoh \ref{con:ruangmetrikkuasi3} biasa disebut ruang metrik kuasi \textit{Sorgenfrey}. Selanjutnya akan diberikan definisi mengenai topologi maju dan topologi mundur.
\begin{defn}\label{def:topologimaju}
    \cite{JuliaCollins2007} Topologi maju $\tau_f$ dari $q$ adalah topologi yang dibangun oleh bola terbuka maju yaitu
    \begin{equation}
        B_f(x,\varepsilon) := \{y\in X\big|q(x,y)<\varepsilon\} \quad \text{untuk }x\in X, \varepsilon>0.
    \end{equation}
    Sedangkan topologi mundur $\tau_b$ dari $q$ adalah topologi yang dibangun oleh bola terbuka mundur yaitu
    \begin{equation}
        B_b(y,\varepsilon) := \{x\in X\big|q(x,y)<\varepsilon\} \quad \text{untuk }y\in X, \varepsilon>0.
    \end{equation}
\end{defn}
Berikut ini, diberikan pula definisi kekonvergenan maju dan mundur pada ruang metrik kuasi.
\begin{defn}\label{def:konvergenmajumundur}
    \cite{JuliaCollins2007} Suatu barisan $(x_n)$ pada suatu ruang metrik kuasi $(X,q)$ disebut 
    \begin{enumerate}
        \item Barisan konvergen maju jika terdapat $x\in X$ sedemikian sehingga 
    \begin{equation*}
        \lim_{n\rightarrow \infty} q(x,x_n) = 0.
    \end{equation*}
    Dengan kalimat lain, untuk setiap $\varepsilon>0$, terdapat $n_0\in \natural$ sedemikian sehingga $q(x,x_n)<\varepsilon$ untuk setiap $n\geq n_0$. Dengan notasi dapat dituliskan $x_n\xrightarrow{f} x$.
        \item Barisan konvergen mundur jika terdapat $x\in X$ sedemikian sehingga 
    \begin{equation*}
        \lim_{n\rightarrow \infty} q(x_n,x) = 0.
    \end{equation*}
    Dengan kalimat lain, untuk setiap $\varepsilon>0$, terdapat $n_0\in \natural$ sedemikian sehingga $q(x_n,x)<\varepsilon$ untuk setiap $n\geq n_0$. Dengan notasi dapat dituliskan $x_n\xrightarrow{b} x$.
    \end{enumerate}
\end{defn}
\begin{defn}\label{def:barisan-cauchy-maju-mundur}
    \cite{JuliaCollins2007} Misalkan $(X,q)$ merupakan ruang metrik kuasi, Barisan $(x_n)$ disebut 
    \begin{enumerate}
        \item Barisan Cauchy maju jika untuk setiap $\varepsilon>0$, terdapat $n_0\in \natural$ sedemikian sehingga berlaku $q(x_n,x_m)<\varepsilon$ untuk setiap $m\geq n\geq n_0$.
        \item Barisan Cauchy mundur jika untuk setiap $\varepsilon>0$, terdapat $n_0\in \natural$ sedemikian sehingga $q(x_m,x_n)<\varepsilon$ untuk setiap $m\geq n\geq n_0$.
    \end{enumerate}
\end{defn}
\begin{defn}\label{def:lengkap-maju-mundur}
    \cite{JuliaCollins2007} Suatu ruang metrik kuasi $(X,q)$ disebut 
    \begin{enumerate}
        \item Ruang metrik kuasi lengkap maju jika setiap barisan Cauchy maju di $X$ konvergen maju di $X$.
        \item Ruang metrik kuasi lengkap mundur jika setiap barisan Cauchy mundur di $X$ konvergen mundur di $X$.
    \end{enumerate}
\end{defn}
Selanjutnya, berikut diberikan pula definisi kekontinuan pada ruang metrik kuasi.
\begin{defn}
    \cite{JuliaCollins2007} Misalkan $(X,q_X)$ dan $(Y,q_Y)$ adalah ruang metrik kuasi. Suatu pemetaan $f:X\rightarrow Y$ dikatakan $f_f$-kontinu, berturut-turut $f_b$-kontinu, pada suatu titik $x_0\in X$ jika untuk setiap $\varepsilon>0$ terdapat $\delta>0$ sedemikian sehingga $y\in B_f(x_0,\delta)$ berakibat $f(y)\in B_f(f(x_0),\varepsilon)$, berturut-turut $f(y)\in B_b(f(x_0),\varepsilon)$.
\end{defn}

Di bawah ini diberikan definisi dari titik tetap pada suatu pemetaan.
\begin{defn}
    \cite{kreyszig} Diberikan $X$ adalah himpunan tak kosong. Titik tetap dari suatu pemetaan $f:X\rightarrow X$ adalah $x\in X$ yang dipetakan ke dirinya sendiri, atau dapat dinotasikan sebagai $f(x)=x$.
\end{defn}
Selanjutnya, berikut ini disajikan hasil titik tetap dengan pemetaan kontraksi$-\psi F$ pada ruang metrik kuasi yang diperoleh Secelean et. al.

Misalkan $\mathcal{F}$ adalah himpunan semua fungsi naik dari $F:(0,+\infty)\rightarrow \R$ dan untuk suatu $\mu\in \R^+,$ $\Psi_{\mu}$ adalah himpunan semua fungsi naik dan kontinu dari $\psi:(-\infty, \mu)\rightarrow\R$ sedemikian sehingga $\psi(t)<t$ untuk setiap $t\in (-\infty,\mu)$. Selanjutnya, $\Psi_{\mu}$ akan dituliskan sebagai $\Psi_F$ jika $F$ suatu fungsi dengan $F\in \mathcal{F}$ dan $\mu\geq \sup F$.

\begin{defn}\label{KontraksipsiF}
    \cite{Secelean2019} Diberikan $(X,q)$ ruang metrik kuasi dan dua fungsi $F\in\mathcal{F}$ serta $\psi\in \Psi_F$. Suatu pemetaan $T:X\rightarrow X$ disebut 
    \begin{enumerate}
        \item[(i)] kontraksi$-\psi F$ maju jika 
                    \begin{align}
                        Tx\neq Ty \implies F(q(Tx,Ty))\leq \psi(F(q(x,y))),  \label{eq:kontraksipsiFkuasimaju}
                    \end{align}
        \item[(ii)] kontraksi$-\psi F$ mundur jika 
                    \begin{align}
                        Tx\neq Ty \implies F(q(Tx,Ty))\leq \psi(F(q(y,x))). \label{eq:kontraksipsiFkuasimundur}
                    \end{align}
    \end{enumerate}
\end{defn}

Berikut ini diberikan definisi dari Operator Picard maju dan mundur yang menjadi poin penting dari teorema titik tetap.
\begin{defn}\label{operatorpicardkuasi}
    \cite{Secelean2019} Suatu pemetaan $T:X\rightarrow X$ pada ruang metrik kuasi $(X,q)$ disebut 
    \begin{enumerate}
        \item[(i)] operator Picard maju jika memiliki suatu titik tetap tunggal yaitu $\zeta\in X$ dan $T^n x\xrightarrow{f} \zeta$ untuk setiap $x\in X$,
        \item[(ii)] operator Picard mundur jika memiliki suatu titik tetap tunggal yaitu $\zeta\in X$ dan $T^n x\xrightarrow{b} \zeta$ untuk setiap $x\in X$.
    \end{enumerate}
\end{defn}

Kemudian, diberikan definisi dari orbit$-T$ dan ruang metrik lengkap maju secara orbit$-T$ yang akan digunakan dalam teorema titik tetap.
\begin{defn}\label{def:orbitT}
    \cite{SHATANAWI2012520} Diberikan $(X,q)$ adalah ruang metrik kuasi dan $T$ adalah pemetaan dari $X$ ke $X$. Himpunan $\mathcal{O}(x,T)=\qty{T^n x~|~n=0,1,2,\cdots}$ disebut orbit dari $T$ pada $x\in X$.
\end{defn}
\begin{defn}\label{def:lengkapmajuorbitT}
    \cite{SHATANAWI2012520} Ruang metrik kuasi $(X,q)$ disebut ruang metrik kuasi lengkap maju secara orbit-$T$ jika setiap barisan Cauchy maju di $\mathcal{O}(x,T), x\in X$, konvergen maju pada suatu titik di $X$.
\end{defn}

Sekarang, diberikan hasil teorema titik tetap dengan kontraksi$-\psi F$ yang diperoleh Secelean et. al. pada \cite{Secelean2019}. 
\begin{teo}\label{teoremaopm}
    \cite{Secelean2019} Diberikan $(X,q)$ ruang metrik kuasi, $F\in \mathcal{F}$, dan $\psi \in \Psi_F$. Jika $T:X\rightarrow X$ adalah pemetaan kontraksi$-\psi F$ mundur dan $X$ ruang metrik kuasi lengkap maju secara orbit$-T$, maka $T$ adalah operator Picard maju.
\end{teo}


%===========================
    \section{HASIL DAN PEMBAHASAN}
    \subsection{Titik Tetap dengan Kontraksi$-\phi G$}
    Dinotasikan $\mathcal{G}$ sebagai himpunan semua fungsi naik tegas dari $G:[0,+\infty)\rightarrow[0,+\infty)$ dan untuk suatu $\mu\in (0,+\infty]$ dinotasikan $\Phi_{\mu}$ sebagai himpunan semua fungsi kontinu naik tegas dari $\phi:[0,\mu)\rightarrow [0,+\infty)$ dengan $\phi(t)<t$ untuk semua $t>0$ dan $\phi(0)=0$. Untuk suatu fungsi $G\in \mathcal{G}$ dan $\mu\geq \sup_{x\in[0,+\infty)} G(x)$, akan dituliskan $\Phi_G$ untuk menyatakan $\Phi_{\mu}$. Berikut ini disajikan contoh dari $G\in G$ dan $\phi\in \Phi_G$.
\begin{con}
    Didefinisikan $G:[0,+\infty)\rightarrow[0,+\infty)$ dengan $G(t):=t^3$, jelas bahwa $G$ naik tegas sehingga $G\in G$. Kita punya $\sup_{x\in[0,+\infty)} G(x)=+\infty$. Selanjutnya, dapat didefinisikan $\phi:[0,\mu)\rightarrow[0,+\infty)$ dengan $\phi(t):=\frac{t}{3}$. Jelas pula bahwa $\phi$ naik tegas dan memenuhi $\phi(0)=0$, serta $\phi(t)=\frac{t}{3}<t$ untuk setiap $t\in [0,\mu)$. Diperoleh $\phi\in \Phi_G$, tetapi $\phi\notin \Psi_G$ karena domain dari $\phi$ adalah $[0,\mu)$ bukan $(-\infty,\mu)$.
\end{con}

Kemudian berikut ini diberikan definisi dari kontraksi$-\phi G$ pada ruang metrik kuasi.
\begin{defn}\label{def:KontraksiPhiG}
Diberikan $(X,q)$ ruang metrik kuasi dan dua fungsi $G\in \mathcal{G}$ serta $\phi\in\Phi_G$. Suatu pemetaan $\Upsilon: X\rightarrow X$ disebut
\begin{enumerate}
    \item[(i)] kontraksi$-\phi G$ maju jika 
    \begin{equation}
        \Upsilon x\neq \Upsilon y \quad \Longrightarrow \quad G(q(\Upsilon x,\Upsilon y)) \leq \phi (G(q(x,y))),
    \end{equation}
    \item[(ii)] kontraksi$-\phi G$ mundur jika
    \begin{equation}
        \Upsilon x \neq \Upsilon y \quad \Longrightarrow  \quad G(q(\Upsilon x, \Upsilon y)) \leq \phi (G(q(y,x))).
    \end{equation}
\end{enumerate}
\end{defn}
Selanjutnya, di bawah ini diberikah hasil teorema titik tetap dengan pemetaan kontraksi$-\phi G$ pada ruang metrik kuasi.
\begin{teo}\label{thm:phiGPicardMaju}
    Misalkan $(X,q)$ ruang metrik kuasi dan dua fungsi $G\in\mathcal{G}$ serta $\phi\in \Phi_{G}$. Jika  $(X,q)$ ruang metrik kuasi lengkap maju yang apabila konvergen maju berimplikasi konvergen mundur, maka untuk setiap kontraksi$-\phi G$ maju $\Upsilon: X\rightarrow X$ adalah operator Picard maju.
\end{teo}

Sebelum membuktikan Teorema \ref{thm:phiGPicardMaju}, diberikan beberapa lemma sebagai pendukung pembuktian sebagai berikut.

\begin{lemmas}\label{lemma:UpsilonKontinu}
    Diberikan $(X,q)$ ruang metrik kuasi dan dua fungsi $G\in\mathcal{G}$ serta $\phi\in \Phi_{G}$. Jika $\Upsilon: X\rightarrow X$ merupakan kontraksi$-\phi G$ maju (berturut-turut mundur), maka 
    \begin{equation}
        q(\Upsilon x,\Upsilon y)\leq q(x,y)\quad \text{berturut-turut } q(\Upsilon x,\Upsilon y)\leq q(y,x)),
    \end{equation}
    untuk setiap $x,y\in X$, sehingga $\Upsilon$ adalah $f_f$ (berturut-turut $f_b$) kontinu.
\end{lemmas}

\begin{proof}
\end{proof}



\begin{lemmas}\label{lemma:phin0}
Jika $\phi\in \Phi_{\mu}$ dengan $\mu>0$, maka $\phi^n(t)\rightarrow 0$ saat $n\rightarrow +\infty$ untuk setiap $t\in[0,\mu)$.
\end{lemmas}
\begin{proof}\label{}
\end{proof}

\begin{lemmas}\label{lemma:Gxnxnterbatas}
    Misalkan $G\in\mathcal{G}$ dan $(x_n)\subseteq [0,+\infty)$. Jika barisan $(G(x_n))$ konvergen dan $\lim_n(G(x_n))<\mu$ dengan $\mu=\sup_{x\in[0,+\infty)}G(x)$, maka barisan $(x_n)$ terbatas.
\end{lemmas}

\begin{proof}
\end{proof}

\begin{lemmas}\label{lemma:Gxnkonvergen0}
     Misalkan $G\in\mathcal{G}$ dan $(x_n)\subseteq [0,+\infty)$. Jika barisan $(G(x_n))$ konvergen ke nol, maka barisan $(x_n)$ juga konvergen ke nol.
\end{lemmas}


\begin{proof}

\end{proof}
\begin{lemmas}\label{lemma:padatR}
    Misalkan $K$ himpunan bagian dari $\R^+$ dengan $K$ padat pada $\R^+$. Jika $x,y\in\R^+$ dengan $x<y$, maka terdapat $\kappa\in K$ sedemikian sehingga $x<\kappa<y$.
\end{lemmas}
\begin{proof}
\end{proof}
\begin{lemmas}\label{3prop}
\cite{Secelean2019} Misalkan $(x_n)$ adalah barisan pada ruang metrik kuasi $(X,q)$ dan $\mathfrak{s}$ adalah himpunan bagian dari $\R^+$ sedemikian sehingga $\R^+\backslash \mathfrak{s}$ padat pada $\R^+$. Jika $q(x_n,x_{n+1})\rightarrow 0$ saat $n\rightarrow+\infty$ dan $(x_n)$ bukan barisan Cauchy maju, maka terdapat $\varepsilon\in \R^+\backslash\mathfrak{s},~n_0\in \natural$ dan barisan bilangan bulat positif $(m_k), (n_k)$ sedemikian sehingga 
\begin{enumerate}
    \item[(a)] $k\leq m_k<n_k$ dan $q(x_{m_k},x_{n_k})\geq \varepsilon$ untuk setiap $k\in\natural$,
    \item[(b)] $n_k\geq m_k+2$ dan $q(x_{m_k},x_{n_k-1})<\varepsilon$ untuk setiap $k\geq n_0$,
    \item[(c)] $q(x_{m_k},x_{n_k})\rightarrow\varepsilon$ saat $k\rightarrow+\infty$.
\end{enumerate}
\end{lemmas}

\begin{proof}

\end{proof}

Akhirnya semua lemma yang diperlukan sudah diberikan. Berikut ini diberikan bukti dari Teorema \ref{thm:phiGPicardMaju}.

\begin{proof}[Bukti dari Teorema]
\end{proof}
\subsection{Aplikasi Teorema Titik Tetap}
Sebelum diberikan aplikasi teorema titik tetap, didefinisikan dulu ruang yang digunakan dan beberapa asumsi yang diperlukan.\\

Misalkan $H$ ruang Banach dan didefinisikan ruang $L^P(0,T;H)$ sebagai berikut.
\begin{equation*}
    L^p(0,T;H)=\{f:[0,T]\rightarrow H:\int_0^T \norm{f(t)}^p_H \, dt <\infty\}, 
\end{equation*}
untuk $p\in[1,+\infty)$, yang dilengkapi oleh norm
\begin{equation*}
    \norm{f}_{L^{p}(0,T;H)} = \left(\int_0^T \norm{f(t)}^p_H \, dt\right)^{\frac{1}{p}}, 
\end{equation*}
untuk $p\in[1,+\infty)$.
Lebih lanjut juga didefinisikan $L^p_{\lambda} (0,T;H)$ sebagai berikut 
\begin{equation*}
    L^p_{\lambda}(0,T;H)=\{f:[0,T]\rightarrow H:\int_0^T e^{-\lambda t}\norm{f(t)}^p_H \, dt <\infty\},
\end{equation*}
untuk $p\in[1,+\infty)$, yang dilengkapi oleh norm 
\begin{equation*}
    \normlp{f} = \left(\int_0^T e^{-\lambda t}\norm{f(t)}^p_H \, dt\right)^{\frac{1}{p}}, 
\end{equation*}
untuk $p\in[1,+\infty)$.

Misalkan $H$ ruang Banach dan $A$ operator linear di $H$. Dalam penelitian ini dibahas mengenai eksistensi dan ketunggalan solusi dari persamaan diferensial elliptik tidak homogen
\begin{equation}
    \dfrac{dw}{dt} = Aw(t) + f(w(t)), \quad w(0)=w_0,~ \label{eq:persdif}
\end{equation}
dengan $f$ fungsi dari $H$ ke $H$. Asumsi-asumsi pada kondisi awal $w_0$, operator $A$, dan fungsi $f$ akan diberikan kemudian setelah diperkenalkan Teorema $\ref{hilleyosida}$. Pertama diperkenalkan beberapa definisi terlebih dahulu. Himpunan resolvent $\sigma(A)$ dari $A$ adalah himpunan semua bilangan kompleks $\lambda$ dengan $\lambda I-A$ mempunyai invers, yakni $(\lambda I -A)^{-1}$ operator linear terbatas di $X$. Untuk $\lambda \in \sigma(A)$, operator linear terbatas $R(\lambda:A):=(\lambda I-A)^{-1}$ disebut resolvent dari $A$. 

\begin{teo}\label{hilleyosida}
    \cite{Pazy1983} Operator linear (tak-terbatas) $A$ adalah infinitesimal generator dari suatu kontraktif $C_0$ semigroup $S(t), t\geq 0$ jika dan hanya jika
    \begin{itemize}
        \item $A$ tertutup\footnote{Operator $A:D(A)\subset X\rightarrow Y$ dikatakan tertutup jika $x_n\rightarrow x$ di $X$ dengan $x_k\in D(A)$ untuk semua $k\in\natural$ dan $Ax_n\rightarrow y$ di $Y$ berimplikasi $x\in D(A)$ dan $y=Ax$.} dan $\overline{D(A)}=H$.
        \item Himpunan resolvent $\sigma(A)$ dari $A$ memuat interval $(0,+\infty)$ dan untuk setiap $\lambda>0$
        \begin{equation*}
            \norm{R(\lambda:A)} \leq \dfrac{1}{\lambda}.
        \end{equation*}
    \end{itemize}
\end{teo}
Dari Teorema di atas, diberikan beberapa asumsi berikut.

\begin{asumsi}\label{asumsi}
\end{asumsi}
Berikut diberikan contoh operator $A$ yang memenuhi kondisi pada Asumsi \ref{asumsi}.

\begin{con}\label{con:asumsiA}
    Misalkan $\Omega$ adalah himpunan bagian dari $\R^n$ yang terbuka dan $A=\Delta$ adalah operator pada $L^2(\Omega)$ yang memenuhi kondisi batas Neumann dan Dirichlet, maka operator $A$ memenuhi Asumsi \ref{asumsi}.
\end{con}
% \begin{proof}
%     Jelas bahwa $D(A)\subset H\rightarrow H$. Ambil sebarang barisan $(x_n)\in D(A)$ yang memenuhi $x_n\rightarrow x$, artinya $\norm{x_n-x}_H\rightarrow 0$. Karena $A$ memenuhi kondisi batas Neumann dan Dirichlet, diperoleh 
%     \begin{align*}
%         0\leq \norm{Ax_n-Ax}_{D(A)} \leq \norm{x_n-x}_H\rightarrow 0.
%     \end{align*}
%     Dari sini didapatkan bahwa $Ax_n\rightarrow y$ dengan $y=Ax$, sehingga $A$ tertutup. Karena $A$ tertutup, diperoleh bahwa $\overline{D(A)}=H$. Selanjutnya, berdasarkan \cite{ricker1988}, kita punya himpunan resolvent $\sigma(A)$ memuat interval $(0,+\infty)$ dan untuk setiap $\lambda>0$, diperoleh
%     \begin{align*}
%         \norm{R(\lambda:A)} \leq \frac{1}{\lambda}\sin^2\qty(\frac{1}{2}\arg(\lambda))\leq \frac{1}{\lambda}.
%     \end{align*}
%     Dengan demikian operator $A$ memenuhi Asumsi \ref{asumsi}.
% \end{proof}

Sebelum dibahas mengenai eksistensi dan ketunggalan solusi dari \eqref{eq:persdif}, terlebih dahulu diperkenalkan definisi solusi \textit{mild} dari \eqref{eq:persdif}. 

\begin{defn}
    \cite{Pazy1983} Misalkan $T>0,p\in [1,+\infty)$ dan diberikan operator $A$ adalah infinitesimal generator dari kontraktif $C_0$ semigroup $S(t), t\geq 0$. Fungsi $u:[0,T]\rightarrow H$ dikatakan solusi \textit{mild} dari \eqref{eq:persdif} apabila ada $\lambda^*=\lambda^* (p,u_0,T)>0$ sehingga $u\in L^p_{\lambda^*}(0,T;H)$ dan memenuhi persamaan 
    \begin{equation}
        u(t)=S(t)u_0+\int_0^t S(t-s)f(u(s))\, ds. \label{eq:solusimild}
    \end{equation}
\end{defn}

Dalam pembuktian eksistensi dan ketunggalan solusi \textit{mild}, terlebih dahulu didefinisikan ruang berikut untuk $R>0$
\begin{align*}
     X^R_{p,\lambda}=\{&g\in L^p_{\lambda}(0,T;H):g(0)=w_0 \\
    &\text{dan } \normlp{g}\leq R\}, \quad \text{untuk } p\in [1,\infty)
\end{align*}
dan didefinisikan fungsi 
\begin{equation*}
    q^R_{p,\lambda}(g,h)=\begin{cases} R, \quad &\text{untuk } h=0 \text{ dan } g\neq 0,\\
    \normlp{g-h}, \quad &\text{untuk } g \text{ dan } h \text{ yang lain}.
    \end{cases}
\end{equation*}
Lebih lanjut didefinisikan pemetaan
\begin{equation}
    \Upsilon(v)(t) = S(t)v_0 + \int_0^t S(t-s)f(v(s)) \, ds. \label{eq:pemetaanT}
\end{equation}

Dalam penelitian ini didapatkan hasil utama sebagai berikut.
\begin{teo}\label{thm:hasilaplikasi}
    Jika $T>0$, $p\in [1,+\infty)$, dan $w_0,~ A, f$ memenuhi Asumsi \ref{asumsi}, maka ada $\lambda^*>0$ dan solusi \textit{mild} yang tunggal dari \eqref{eq:persdif} di $X^R_{p,\lambda^*}$.
\end{teo}

Sebelum memberikan bukti Teorema \ref{thm:hasilaplikasi} diberikan beberapa lemma berikut sebagai pendukung pembuktian.
\begin{lemmas}
    Misalkan $p\in[1,+\infty)$ dan $\lambda, R>0$. Ruang $(X^R_{p,\lambda},q^R_{p,\lambda})$ adalah ruang metrik kuasi.
\end{lemmas}

\begin{proof}
\end{proof}

\begin{lemmas}\label{lemma:X_metrik_kuasi_lengkapmaju}
    Misalkan $p\in[1,+\infty)$ dan $\lambda, R>0$. Ruang $(X^R_{p,\lambda},q^R_{p,\lambda})$ adalah ruang metrik kuasi lengkap maju.
\end{lemmas}

\begin{proof}

\end{proof}

\begin{lemmas}\label{lemma:konvergen_maju-konvergen_mundur}
    Misalkan $p\in[1,+\infty)$ dan $\lambda, R>0$. Ruang $(X^R_{p,\lambda},q^R_{p,\lambda})$ adalah ruang metrik kuasi yang apabila konvergen maju, maka konvergen mundur.
\end{lemmas}

\begin{proof}

\end{proof}

\begin{lemmas}\label{lemma:UpsilonXkeX}
    Misalkan $p\in[1,+\infty)$ dan $\lambda, R>0$. Jika $w_0,~A,$ dan $f$ memenuhi Asumsi \ref{asumsi}, maka $\Upsilon$ memetakan dari $X^R_{p,\lambda^*}$ ke dirinya sendiri dengan 
    \begin{align}
        \nonumber \lambda^*\geq &~\dfrac{(2^{p-1}+2^{3p-3}K^pT^p)\norm{v_0}^p_H +2^{2p-2}T^p\norm{f(v_0)}^p_H}{R^p} \\
        &+ 2^{3p-3}K^pT^{p-1}.\label{eq:ketaksamaanLambda1}
    \end{align}
\end{lemmas}

\begin{proof}

\end{proof}

\begin{lemmas}\label{lemma:qTu<qu}
    Misalkan $p\in [1,+\infty)$ dan $R>0$. Jika $w_0,~A,$ dan $f$ memenuhi Asumsi \ref{asumsi}, maka 
    \begin{equation}
        \qpl(\Upsilon u, \Upsilon v)\leq \alpha \qpl(u,v), 
    \end{equation}
    untuk semua $u,v\in X^R_{p,\lambda^*}$ dengan $\lambda^*>T^{p-1}K^p$ dan $\alpha = \left(\dfrac{T^{p-1}K^p}{\lambda^*}\right)^{\frac{1}{p}}$.
\end{lemmas}

\begin{proof}

\end{proof}
Akhirnya semua lemma yang diperlukan sudah diberikan. Di bawah ini, diberikan bukti dari Teorema \ref{thm:hasilaplikasi}.

\begin{proof}[Bukti dari Teorema \ref{thm:hasilaplikasi}]

\end{proof}




%==================================================
    \section{KESIMPULAN}
    Tuliskan kesimpulan dari penelitian yang artikelnya Anda tulis ini tanpa mengulang hal-hal yang telah disampaikan di Abstrak. Kesimpulan dapat diisi pula tentang pentingnya hasil yang dicapai dan saran untuk aplikasi dan pengembangannya.
    \section*{LAMPIRAN}
    Jika ada, lampiran muncul di sini.
    \section*{UCAPAN TERIMA KASIH}
    Tuliskan ucapan terima kasih dengan bahasa baku, misalnya, Penulis A.F. mengucapkan terima kasih kepada Direktorat Pendidikan Tinggi, Departemen Pendidikan dan Kebudayaan Republik Indonesia yang telah memberikan dukungan finansial melalui Beasiswa Bidik Misi tahun 2010-2014. Penulis juga diperkenankan menyampaikan ucapan terima kasih kepada sponsor penyedia dana penelitian.
    
    \section*{DAFTAR PUSTAKA}
    \bibliographystyle{IEEEtran}
    \bibliography{file/daftarPustaka}
       
\end{document}